\documentclass[11pt,a4paper]{article}
\usepackage[T1]{fontenc}
\usepackage[utf8]{inputenc}
\usepackage[margin=2.5cm]{geometry}
\usepackage{amsmath,booktabs,tabularx,array}
\usepackage[hidelinks]{hyperref}
\newcommand{\code}[1]{\texttt{\detokenize{#1}}}
\newcolumntype{Y}{>{\raggedright\arraybackslash}X}
\setlength{\parindent}{0pt}
\setlength{\parskip}{0.6em}

\title{Wearable Pipeline: Data, Training, Inference and Overlap}
\author{}
\date{Code-reading notes: 4 November 2026}

\begin{document}
\maketitle

This describes the current implementation, rather than assuming every statement in
the older documentation remains current. No project code or tests were run for this
review. The concerns below come from static inspection and have not been reproduced
in a fresh experiment.

\section{The main distinction}

The pipeline has two kinds of learning: estimating each patient's signal behaviour,
and training a classifier to recognise exacerbations. They happen at different stages.

\begin{center}
Raw recordings and clinical labels\\
$\downarrow$\\
Separate channel files and a daily label table\\
$\downarrow$\\
Raw / smoothing / frozen forecasting / frozen encoding\\
$\downarrow$\\
Optional change-point detection\\
$\downarrow$\\
Daily features and history windows\\
$\downarrow$\\
Patient-level train/test split\\
$\downarrow$\\
Supervised classifier $\longrightarrow$ held-out probabilities and evaluation
\end{center}

The signal-processing routes are alternative experimental branches. They do not all
have to run consecutively.

\section{Raw data to channel files and labels}

The default source reads \code{data/raw/pass1/patients.csv}, selects patients marked
\code{in_models}, and loads their minute-level recordings.

\begin{itemize}
\item Clinical labels are collapsed to one row per patient-day before filtering
wearable readings. The labels already exist in the export; this repository reads them.
\item HRV, heart rate, temperature, steps, sleep and activity channels are extracted
separately. Each channel retains its own observation times.
\item Long gaps receive NaN filler rows. These mark missingness without inventing
readings. Filtering chunks are defined from actual observations, with a gap of at
least 180 minutes opening a new chunk.
\end{itemize}

Nothing is trained here. Output is one processed file per patient-channel and a
separate daily ground-truth table, \code{data/processed/labels_pass1.csv}.

\textbf{Code:} \code{wearable/sources.py}; \code{wearable/steps/ingest.py};
\code{wearable/chunking.py}.

\section{What each signal-processing route learns}

\begin{tabularx}{\textwidth}{@{}p{0.20\textwidth}YY@{}}
\toprule
Route & What is fitted or trained? & Output \\
\midrule
Raw & Nothing. & Original readings. \\
Smoothing & Patient-specific parameters and evolving filter states; no clinical
labels. & Baseline and residual: raw minus baseline. \\
Frozen forecasting & Timestral, TimesFM or Chronos pretrained weights remain frozen.
& Predicted readings, forecast error (surprise), and distribution position (tail). \\
Frozen encoding & MOMENT or Mantis pretrained weights remain frozen. Each input
window is normalised separately. & One representation vector per patient-day. \\
\bottomrule
\end{tabularx}

\textbf{The frozen forecasters predict physiological readings, not exacerbations.}
Their outputs become features for a separate clinical classifier. Forecasting uses
preceding observations to predict blocks ahead. Encoding instead reads a window
including the day being classified and summarises it as a vector.

The experiment normally processes HRV and HR separately. Other available channels
join later as raw features. Chronos is currently used as a univariate forecaster;
its possible cross-channel covariate route is not implemented.

\textbf{Code:} \code{wearable/experiment.py} (\code{inputs});
\code{wearable/steps/smooth.py}; \code{wearable/steps/forecast.py};
\code{wearable/steps/encode.py}.

\section{Optional change-point detection}

Change-point detection (CPD) can read raw values, smoothed values, smoothing residuals,
forecast surprise, or forecast tail. For each patient and channel it:

\begin{enumerate}
\item Estimates an hour-of-day and behavioural-state baseline from prior readings.
\item Converts deviations into standardised innovations using a noise model.
\item Calibrates detector thresholds using synthetic no-change histories built from
that patient's past.
\item Applies BOCPD, KCPD, Kalman, HMM and nightly drift detection.
\item Combines detector alarms using OR.
\end{enumerate}

This is patient-specific, unsupervised adaptation, including on test patients' own
histories. It does not use exacerbation labels. Default calibration shuffles historical
days; optional STS calibration fits a generative time-series model to past readings.
The noise model refits from trailing history, and alarm thresholds are recalibrated
periodically.

The output contains flags, magnitudes, noise statistics and gap-reentry information.
It is not a probability of exacerbation. Causal detector processing also cannot undo
future leakage already introduced by its input smoother.

\textbf{Code:} \code{wearable/steps/detect.py}; \code{wearable/cpd/_whiten.py};
\code{wearable/cpd/_sts.py}.

\section{Readings to classifier examples}

The \code{join} step builds one example per patient-day:

\begin{itemize}
\item Continuous readings become daily mean, standard deviation, minimum and maximum.
\item State/activity and alarm columns contribute totals and bout counts.
\item Optional profiles preserve time-of-day bins instead of only daily summaries.
\item A seven-day input window attaches day $D$ and days $D-1$ through $D-6$.
Missing calendar days remain missing; older days do not shift into their places.
\item Embeddings already represent the requested window, so they are not lagged again.
\item Clinical labels are attached by patient and date. Identifiers, clinical labels,
event IDs and cohort group are excluded from classifier features.
\end{itemize}

\begin{tabularx}{\textwidth}{@{}p{0.15\textwidth}YY@{}}
\toprule
Task & Available features & Target \\
\midrule
Detection & Through day $D$, including $D$. & Is $D$ an exacerbation day? \\
Onset & Through day $D$, including $D$. & Does a new episode start within the next
$H$ days? \\
\bottomrule
\end{tabularx}

The tasks produce separate datasets and results. The default onset horizon is three
days; it can be narrowed to one or two while retaining the same eligible days.

\textbf{These are full-day predictions.} Using all readings from $D$ supports a
decision after those readings are available. It does not establish an alert available
at the beginning of $D$.

The default eligibility rule retains the study's HRV-based sample days even when
other channels are included as features. An optional wider validity rule exists;
changing eligibility changes the population and needs separate review.

\textbf{Code:} \code{wearable/steps/join.py}; \code{wearable/windowing.py}.

\section{Supervised training and held-out inference}

Patients are split 80:20 by default, stratified by cohort and whether they have an
exacerbation. All days of a patient stay on one side. Five seeds repeat the experiment.
Within a comparison, input methods use the same split.

The final model receives training features and training labels:

\begin{itemize}
\item \textbf{XGBoost}, the default, learns boosted trees; CatBoost is another tree
classifier. Both handle class imbalance during training.
\item \textbf{RNN/LSTM} learn sequences of daily feature vectors. Feature scales are
estimated from training data and reused for test inference.
\item \textbf{Mantis classifier} uses pretrained Mantis and trains it for the clinical
task. The default requests full fine-tuning, but when a channel adapter is required,
the implementation switches to adapter/head training.
\item \textbf{Always-no} is the reference that never flags a day.
\end{itemize}

The fitted model then predicts $P(y=1)$ for held-out patients. It does not receive
their clinical labels during prediction. Evaluation subsequently compares probabilities
with those labels, computes metrics, resamples test patients for confidence intervals,
and writes probability-band calibration tables.

\textbf{The calibration table assesses probabilities; it does not recalibrate them.}
It is also distinct from the unsupervised alarm-threshold calibration in CPD.

The model step saves predictions and provenance, not a fitted classifier checkpoint.
It is currently an experiment workflow rather than a separate saved-model deployment
flow. The artifact store caches results so the same signal-processing outputs can be
reused across tasks, models and folds.

\textbf{Code:} \code{wearable/steps/model.py}; \code{wearable/models/};
\code{wearable/steps/evaluate.py}; \code{wearable/store.py}.

\section{Where the data and features overlap}

\begin{enumerate}
\item \textbf{Adjacent days.} Seven-day windows overlap by six days within a patient.
This is expected. Patient-level splitting prevents these windows from crossing
train/test within a run.
\item \textbf{Different seeds.} A patient may train in one run and test in another.
These are repeated holdouts, not five independent cohorts.
\item \textbf{Different branches.} A smoothed artifact carries raw values, baseline
and residual. A residual-plus-CPD route passes the smoothing block and detector
features. An embedding route adds embeddings beside ordinary raw features.
Therefore these comparisons concern feature bundles, not isolated representations.
\item \textbf{Before the split.} Signal processing runs before the patient split and
is reused across folds. Frozen inference and adaptation using a test patient's
available past can be legitimate. Using that patient's future is not causal.
\item \textbf{Signal-only smoothing studies.} The separate smoothing-study script
can derive settings from entire cohorts. If those settings are later used for a
held-out predictive claim, their derivation must also respect the training boundary;
being label-free alone does not establish an untouched test set.
\end{enumerate}

\section{Concrete concerns to discuss before corrections}

\subsection{Future-dependent smoother calibration}
Nine of the ten smoothers estimate patient-specific parameters from the full supplied
recording. Later observations can therefore change earlier features. Patient-level
splitting does not fix this within-patient lookahead. The documented truncation-exact
exception is \code{anrewma}; this review did not rerun that check.

For example, \code{gammadglm} fits its patient calibration before filtering individual
chunks. A forward filtering loop does not make the whole operation causal when its
parameters were estimated using future readings.

\subsection{CPD channel bookkeeping}
Detect artifacts do not record \code{stats.signal}, but \code{join._features} uses
that field to decide which raw channel an artifact already covers. The fallback is
HRV. An HR detector can therefore be treated as covering HRV, changing which raw
feature blocks are included. This needs correction or explicit verification before
assuming the intended variables reach every path.

\subsection{Flag counts labelled as minutes}
Daily aggregation sums flagged rows into columns ending in \code{_minutes}. For
irregular HRV readings, a flagged-row count is not elapsed time. The feature should
either be named as a count or computed using an explicit duration rule.

\section{Scope of these notes}

These notes explain the implemented flow and identify specific points for discussion.
They do not certify causality, vendor-model behaviour, clinical performance, or existing
cached results. No changes to pipeline code were made, and no experiments, tests or
LaTeX compilation were run for this document.

\end{document}
